% Part: counterfactuals
% Chapter: minimal-change-semantics
% Section: transitivity

\documentclass[../../../include/open-logic-section]{subfiles}

\begin{document}

\olfileid{cnt}{min}{tra}

\olsection{Transitivitas}

Kanggo kondisional material, aturan rantai lumaku: $!A \lif !B, !B
\lif !C \Entails !A \lif !C$. Kanthi tembung liya, kondisional
material transitif. Apa iki uga lumaku kanggo kontrafaktual?
Gatekna tuladha saka Stalnaker iki.
\begin{quote}
  Saupama J.~Edgar Hoover lair minangka wong Rusia, dheweke mesthi dadi Komunis.

  Saupama J.~Edgar Hoover dadi Komunis, dheweke mesthi dadi pengkhianat.

  Mula, saupama J.~Edgar Hoover lair minangka wong Rusia,
  dheweke mesthi dadi pengkhianat.
\end{quote}
Saupama Hoover lair (ing wektu kang padha karo laire kang nyata)
ing Rusia, dudu ing Amerika Serikat, dheweke mesthi gedhe ing
Uni Soviet lan dadi Komunis (upamane kita nganggep mangkono).
Mula premis kapisan bener. Premis kapindho, yen ditimbang dhewe,
uga bener. Nanging konklusine luput: kemungkinan gedhe Hoover
mesthi dadi Komunis kang temen saupama lair ing USSR, lan ora
dadi pengkhianat marang negarane. Penetapan nilai bebener
kanthi intuitif iki diterangake dening panjelasan Stalnaker--Lewis.
Donya kang bisa ana lan paling cedhak karo donya kita, kanthi
owah-owahan mung panggonan laire Hoover, yaiku donya kang ing
kono Hoover gedhe dadi warga USSR kang becik. Iki donya paling
cedhak kang nyukupi anteseden premis kapisan lan konklusi.
Ing donya iku Hoover anggota Partai Komunis kang setya, mula
dudu pengkhianat. Nanging kanggo ngevaluasi premis kapindho,
kita kudu nimbang donya kang beda: donya paling cedhak kang
ing kono Hoover dadi Komunis, yaiku donya kang ing kono dheweke
lair ing Amerika Serikat, ganti kasetyan, lan mula dadi
pengkhianat.\footnote{Mesthi wae, kanggo mangerteni daya tuladha
  iki kita kudu nampa sawetara asumsi metafisika lan politik,
  upamane yen Hoover bisa wae lair saka wong tuwa Rusia, utawa
  yen para Komunis ing Amerika Serikat taun 1950-an iku
  pengkhianat marang negarane.}

\begin{prob}
  Golekana tuladha intuitif kang nggawe yakin tumrap gagale
  transitivitas kontrafaktual.
\end{prob}

\begin{ex}\ollabel{ex:trans-counterex}
  Semantik bola nuduhake yen inferensi iki ora sah, yaiku $p
  \cif q, q \cif r \Entails/ p \cif r$. Gatekna model
  $\mModel{M} = \tuple{W, O, V}$ kanthi $W = \{w, w_1, w_2\}$,
  $O_w = \{\{w\}, \{w, w_1\}, \{w, w_1, w_2\}\}$,
  $V(p) = \{w_2\}$, $V(q) = \{w_1, w_2\}$, lan $V(r) = \{w_1\}$.
  Ana bola kang ngakoni $p$, yaiku $S = \{w, w_1, w_2\}$, lan
  $p \lif q$ bener ing kabeh donya ing kono, mula
  $\mSat{M}{p \cif q}[w]$. Ana uga bola kang ngakoni $q$, yaiku
  $S' = \{w, w_1\}$, lan $q \lif r$ bener ing kabeh donya ing
  kono, mula $\mSat{M}{q \cif r}[w]$. Nanging bola siji-sijine
  kang ngakoni $p$, yaiku $\{w, w_1, w_2\}$, ngemot donya~$w_2$
  kang ing kono $\mSat/{M}{p \lif r}[w_2]$. Mula konklusi
  kontrafaktual luput ing donya pusat.
\end{ex}

\begin{prob}
Gambarna diagram bola kang cocog karo tuladha panyanggah ing
\olref[cnt][min][tra]{ex:trans-counterex}.
\end{prob}

\begin{prob}
  Ing \olref[cnt][min][tra]{ex:trans-counterex}, donya $w_2$ yaiku
  donya kang ing kono Hoover lair ing Rusia, dadi komunis lan
  dudu pengkhianat; $w_1$ yaiku donya kang ing kono Hoover lair
  ing Amerika Serikat, dadi komunis lan dadi pengkhianat.
  Ing model iki, $w_1$ luwih cedhak karo~$w$ tinimbang $w_2$.
  Apa iki perlu? Apa kita bisa menehi tuladha panyanggah kang
  ora nganggep kemungkinan Hoover lair ing Rusia luwih adoh
  tinimbang kemungkinan dheweke dadi Komunis?
\end{prob}

\end{document}
